%%%=========[BT_01]================%%%
\begin{bt}%[KHTN6]
Hãy xác định trạng thái của các hỗn hợp sau đây bằng cách điền vào bảng: với mỗi hỗn hợp, đánh dấu \textbf{x} vào cột \textbf{Dung dịch}, \textbf{Huyền phù} hoặc \textbf{Nhũ tương} cho phù hợp.

\begin{center}
\begin{tabular}{|l|c|c|c|}
\hline
\textbf{Hỗn hợp} & \textbf{Dung dịch} & \textbf{Huyền phù} & \textbf{Nhũ tương} \\
\hline
Sữa chua lên men       & & & \\
\hline
Hoà đất vào nước       & & & \\
\hline
Hoà muối ăn vào nước   & & & \\
\hline
Hoà đường vào nước     & & & \\
\hline
Sữa tươi               & & & \\
\hline
Dầu gội đầu            & & & \\
\hline
Sữa tắm                & & & \\
\hline
\end{tabular}
\end{center}

\loigiai{
\textbf{Phân tích từng hỗn hợp:}

\begin{itemize}
    \item \textbf{Dung dịch} là hỗn hợp đồng nhất, trong suốt, chất tan tan hoàn toàn vào dung môi (thường là nước). Ví dụ: muối, đường tan trong nước.
    \item \textbf{Huyền phù} là hỗn hợp không đồng nhất, gồm các hạt rắn lơ lửng trong chất lỏng, thường đục và có thể lắng xuống theo thời gian.
    \item \textbf{Nhũ tương} là hỗn hợp không đồng nhất, gồm các giọt chất lỏng lơ lửng trong một chất lỏng khác (hai chất lỏng không tan vào nhau). Thường có dạng sệt, mịn.
\end{itemize}

\textbf{Bảng đáp án:}

\begin{center}
\begin{tabular}{|l|c|c|c|}
\hline
\textbf{Hỗn hợp} & \textbf{Dung dịch} & \textbf{Huyền phù} & \textbf{Nhũ tương} \\
\hline
Sữa chua lên men       &   &   & x \\
\hline
Hoà đất vào nước       &   & x &   \\
\hline
Hoà muối ăn vào nước   & x &   &   \\
\hline
Hoà đường vào nước     & x &   &   \\
\hline
Sữa tươi               &   &   & x \\
\hline
Dầu gội đầu            &   &   & x \\
\hline
Sữa tắm                &   &   & x \\
\hline
\end{tabular}
\end{center}

\textbf{Giải thích chi tiết từng trường hợp:}

\begin{enumerate}
    \item \textbf{Sữa chua lên men:} Sữa chua là hỗn hợp gồm các giọt chất béo (lipid) phân tán trong pha nước chứa protein và acid lactic. Đây là \textbf{nhũ tương}.

    \item \textbf{Hoà đất vào nước:} Các hạt đất (khoáng chất, sét) không tan trong nước mà lơ lửng dưới dạng hạt rắn nhỏ. Đây là \textbf{huyền phù}. Để lâu, đất sẽ lắng xuống đáy.

    \item \textbf{Hoà muối ăn vào nước:} Muối ăn ($NaCl$) tan hoàn toàn trong nước, tạo thành hỗn hợp đồng nhất, trong suốt. Đây là \textbf{dung dịch}.

    \item \textbf{Hoà đường vào nước:} Đường ($C_{12}H_{22}O_{11}$) tan hoàn toàn trong nước, tạo hỗn hợp đồng nhất, trong suốt. Đây là \textbf{dung dịch}.

    \item \textbf{Sữa tươi:} Sữa tươi gồm các giọt chất béo phân tán trong pha nước chứa protein, đường lactose và khoáng chất. Đây là \textbf{nhũ tương} (fat-in-water emulsion).

    \item \textbf{Dầu gội đầu:} Dầu gội gồm các thành phần dầu, nước và chất hoạt động bề mặt (surfactant) tạo thành hệ nhũ hoá bền vững. Đây là \textbf{nhũ tương}.

    \item \textbf{Sữa tắm:} Tương tự dầu gội, sữa tắm là hỗn hợp nhũ hoá giữa dầu và nước nhờ chất hoạt động bề mặt. Đây là \textbf{nhũ tương}.
\end{enumerate}

\textbf{Lưu ý sư phạm:}
Khi phân biệt ba loại hỗn hợp này, học sinh cần nhớ:
\begin{itemize}
    \item Nhìn bằng mắt thường: dung dịch \textbf{trong suốt}, huyền phù và nhũ tương \textbf{đục}.
    \item Để yên: huyền phù \textbf{lắng xuống}, nhũ tương \textbf{không lắng} (hoặc lắng rất chậm).
    \item Thành phần: huyền phù chứa \textbf{chất rắn} lơ lửng; nhũ tương chứa \textbf{giọt lỏng} lơ lửng.
\end{itemize}
}
\end{bt}
